\chapter{Conclusions and Future Works}
\label{chap:conclusions}

\section{Conclusions}

This dissertation addressed a practical question in the monitoring of mobile networks: at which level of aggregation should telecommunications traffic be analysed to detect anomalies?
A fully aggregated series hides anomalies that affect a small segment of the traffic, while one series per combination of categories produces too many sparse series to monitor and interpret.
Instead of fixing the aggregation in advance, this work treated it as a combinatorial search.
Candidate hierarchies were generated by a \ac{BNF} grammar over the categorical features of the traffic, evolved with \ac{GE} in PonyGE2, and ranked by how well a label-free Matrix Profile detector found synthetic anomalies in the series they produced.

The framework was applied to hourly \#Events from a mobile network switch, covering July to October 2022, with about one million records and six categorical features.
Because the operational data has no reliable anomaly labels, 105 point anomalies were injected at three standard deviations from the local weekly mean, each in a randomly chosen segment of the traffic.
The data analysis showed that the traffic is highly concentrated: a few combinations carry most of the events, while most of the 1{,}555 non-empty combinations are small and sparse.
Partner Description, with 392 values and the highest entropy, was therefore excluded from the grammar, which kept the search space manageable and avoided splits that produce hundreds of nearly empty series.

The main conclusions are the following.

\begin{itemize}
    \item \textbf{Aggregation matters as much as the detector.}
    With the detector held fixed, the choice of aggregation alone changed the test \ac{AUC} from 0.54 for the fully aggregated series to 0.88 for full disaggregation.
    The injection experiments also showed why: an anomaly that is obvious in its own segment is diluted, and often invisible, in any series that also contains other traffic.
    This supports the premise of the dissertation that the view of the data should be chosen, and not assumed.

    \item \textbf{A grammar is a practical way to define the search space.}
    The grammar generates only valid hierarchies, each leaf is a readable conjunction of business filters, and the leaves always partition the traffic, so no repair or penalty is needed.
    The space it defines contains about $2.2 \times 10^{61}$ distinct hierarchies, which rules out exhaustive search, yet PonyGE2 found hierarchies with 8 to 34 leaves in 510 evaluations per run.

    \item \textbf{The evolved hierarchies are competitive, but not the best overall.}
    They reached a test \ac{AUC} of $0.678 \pm 0.066$ and were significantly better than the fully aggregated series and the Call, Event and Service Type splits.
    They did not outperform the best single-feature splits, and under the original metric full disaggregation was better.
    Their advantage is a compact, interpretable set of about 20 series that is found without knowing in advance which features to split on.

    \item \textbf{The evaluation metric changes the conclusions.}
    The event-weighted \ac{AUC} used as fitness ignores leaves without anomalies, so false alarms on those leaves are not counted.
    When they are counted, with a pooled \ac{AUC}, full disaggregation falls from 0.877 to 0.600, with a precision of 0.4\%, while the evolved hierarchies keep a similar value (0.670).
    The two runs with the best original fitness had the lowest pooled \ac{AUC}, which shows that the search can exploit this blind spot of the fitness function.
    This is the main methodological finding of the work.

    \item \textbf{The Matrix Profile works well on dense series and fails on sparse ones.}
    On high-volume series with a regular weekly shape, the first difference of the profile located the injected anomalies with leaf \ac{AUC}s up to 1.00.
    On sparse series with a few events per hour, such as inbound roaming segments, every isolated spike looks like a discord and the detector performed worse than random.
    Which leaves a hierarchy creates therefore determines not only how visible the anomalies are but also whether the detector can work at all.

    \item \textbf{Training fitness was a weak guide to test performance.}
    Most runs lost 0.10 to 0.15 \ac{AUC} on the test period, and the fitness landscape was flat, with very different hierarchies reaching similar fitness.
    The small search budget, ten individuals for 50 generations, also meant that the search had not converged.
\end{itemize}

The contributions of this dissertation are, first, the formulation of aggregation choice for anomaly detection in grouped telecommunications series as a grammar-based search, with a fitness built from a label-free detector and synthetic anomalies; second, an implementation in Python on top of PonyGE2 and STUMPY that evaluates candidate hierarchies on real operator data and can be reproduced from the saved runs; and third, an empirical analysis that compares evolved and fixed hierarchies under two metrics and shows that conclusions about the best aggregation depend strongly on how false alarms are counted.
For a network operations team, the practical message is that monitoring at an intermediate level of aggregation, chosen with the data, gives a small number of interpretable series in which anomalies are more visible than in the total, as long as sparse segments are avoided.

\section{Future Work}

The results and limitations of this work suggest several directions.

\begin{itemize}
    \item \textbf{A fitness that counts false alarms.}
    The most direct improvement is to use the pooled \ac{AUC}, or another metric that includes anomaly-free leaves, as the fitness.
    It is cheap to compute and would remove the incentive, observed in the best runs, to create many normal-only leaves.
    Precision-oriented metrics, such as the precision–recall \ac{AUC} or F1 at an operational alarm budget, are also worth testing.

    \item \textbf{Control of sparse leaves.}
    Since the Matrix Profile fails on sparse series, a minimum-support rule could stop a branch when its traffic falls below a threshold, either in the grammar or in the fitness function.
    A multi-objective formulation could additionally trade detection quality against the number of leaves, making the cost of monitoring an explicit objective.

    \item \textbf{A larger search budget and a study of stability.}
    With more individuals and generations, it would be possible to establish whether the search converges, how sensitive it is to the random seed, and whether the flat fitness landscape contains a stable core of useful splits.

    \item \textbf{Real incidents and other anomaly types.}
    Only synthetic point anomalies were evaluated.
    Validating the evolved hierarchies against real incident records from the operator, and injecting collective or contextual anomalies that last several hours, would show whether the conclusions hold for the anomalies that matter in operation.

    \item \textbf{More robust validation over time.}
    A single temporal split was used.
    Rolling-origin evaluation over several periods, and data from months beyond July to October 2022, would give a better estimate of how well an evolved hierarchy generalises and how often it needs to be re-evolved.

    \item \textbf{Reintroducing Partner Description.}
    The partner feature was excluded because of its cardinality.
    Grouping partners, for example keeping the most frequent ones and merging the long tail into a single category, would let the search use this information without creating hundreds of sparse branches.

    \item \textbf{Other scores and detectors.}
    Only the first difference of the Matrix Profile was tested.
    The raw profile and alternative detectors, such as forecast residuals or autoencoders, could be plugged into the same fitness function, keeping \ac{GE} for the structure search and testing whether the best aggregation depends on the detector.

    \item \textbf{Additional targets.}
    Only \#Events was used.
    Event Duration could be added as a second target, either as a separate search or as a second objective, since some faults may affect call length without changing the number of calls.

    \item \textbf{Towards operation.}
    Finally, the approach could be moved towards an online setting, in which the hierarchy is evolved periodically on recent data and the leaf series are scored as new hours arrive, with the evolved tree and the discord plots presented to operators as an explanation of each alert.
\end{itemize}
